\documentclass[11pt]{article}
\usepackage[a4paper,margin=24mm]{geometry}
\usepackage[T1]{fontenc}
\usepackage{lmodern,amsmath,amssymb,amsthm,mathtools,booktabs,microtype}
\usepackage[colorlinks=true,linkcolor=blue,citecolor=blue,urlcolor=blue]{hyperref}
\newcommand{\eps}{\varepsilon}
\newcommand{\sech}{\operatorname{sech}}
\newcommand{\tr}{\tau}
\newtheorem{theorem}{Theorem}
\newtheorem{lemma}{Lemma}
\title{Preparation memory in the neutral slope\\
\large An exact two-step Mehler kernel and its string-time limit}
\author{Alan Day \and Chengheng\thanks{AI research assistant (OpenAI).}}
\date{8 November 2026\quad Local research draft}
\begin{document}\maketitle
\begin{abstract}
The first neutral coefficient of the actual bilocal retains
preparation time even when its charged curvature restriction is
stationary. We reduce its selected two-time function exactly to an
energy-tilted two-point function plus a positive Gram kernel formed
by two q-Mehler transitions. A weighted three-energy estimate gives
the compact-time small-deformation expansion with an
$O(\varepsilon^2)$ remainder. The additional term is separable at
first order and survives the equal mixture that cancels the
curvature correction. The complementary odd mixed source produces
an order-$\sqrt{\varepsilon}$ energy--slope expectation.
This establishes a concrete neutral
contribution to the mixed source, not a complete bilocal or bulk
particle identification.
\end{abstract}

\section{Objects, actual interface and result}
Let $h,U$ be orthogonal unit-variance q-Gaussian fields in the
tracial covariant algebra, and put
\[
 q=e^{-\eps},\quad \lambda=\eps/4,\quad s=\sqrt q,\quad
 W=\sqrt\lambda U,\quad K=h/\sqrt\eps,\quad U(u)=e^{iuK}Ue^{-iuK}.
\]
The dimensionless string time is $u=2J_st$, with $\hbar=1$.
The already established single-selected density \cite{selected} is
\begin{equation}\label{density}
 \rho=1-s(1+\lambda)+s(h-\sqrt\lambda U)^2,\qquad
 \sigma(Y)=\tr(\rho Y).
\end{equation}
The normalized actual exact-pinched neutral slope converges to U
inside these preparations \cite{actual}. Thus the functions below
refer to a controlled microscopic coefficient. Take finite-N exact
projection, then half-filled fixed-$\lambda$ double scaling with
the original common disorder, then $\eps\downarrow0$.

Write $g(u)=\sech^2u$ and
$R(u)=\tfrac12g(u)[u\tanh u-3\tanh^2u]$.
\begin{theorem}\label{main}
For fixed $S,k<\infty$, uniformly on $|u|,|v|\le S$ including
all partial derivatives of total order at most k,
\begin{align}
 \tr(U(u)U(v))&=g(u-v)-\tfrac{\eps}{2}R(u-v)+O(\eps^2),\nonumber\\
 \sigma(U(u)U(v))&=g(u-v)+\tfrac{\eps}{2}
       [R(u-v)+g(u)g(v)]+O(\eps^2).\label{result}
\end{align}
\end{theorem}
U has unit \emph{tracial} variance. Its selected variance at the
preparation time is $1+s\lambda(1+q)$, not one.
There is no N-uniform rate or growing-time limit in this statement.

\newpage
\section{Four fields become a classical three-energy chain}
Let $\mu_q$ be the unit-variance q-Hermite law and $P_r$ the
q-Mehler operator,
\[
 P_rH_n=r^nH_n,\qquad
 P_rf(a)=\int K_r(a,b)f(b)\,\mu_q(db).
\]
Its positivity, trace preservation and semigroup law are standard
\cite{bks}. Define the probability measure
\begin{equation}\label{chainmeasure}
 d\Pi_q(a,b,c)=
 \mu_q(da)K_q(a,b)\mu_q(db)K_q(b,c)\mu_q(dc).
\end{equation}
The subscript q on $K_q$ denotes its transition parameter, with
the same deformation parameter q understood.

\begin{lemma}\label{four}
For bounded continuous f,g,k on the energy spectrum,
\begin{equation}\label{fouridentity}
 \tr[(U^2-1)f(h)U g(h)U k(h)]
   =(1+q)\int f(a)g(b)k(c)\,d\Pi_q.
\end{equation}
\end{lemma}
\begin{proof}
It suffices first to use monomials. Label the four U endpoints
in cyclic order by $0,1,2,3$, with no h's between 0 and 1.
The pairing $(0,1)(2,3)$ is removed by $U^2-1$.
The two remaining pairings are $(0,3)(1,2)$ and
$(0,2)(1,3)$, with mutual U-crossing weights 1 and q.
In both cases an h-chord joining the f and g blocks crosses one
U-chord, one joining g and k crosses one, and one joining f and k
crosses two. Within-block weights are unchanged. After summing
the h-pairings, both contributions have the q-Gaussian three-block
law with covariance
\[
 \begin{pmatrix}1&q&q^2\\q&1&q\\q^2&q&1\end{pmatrix}.
\]
By the known q-Gaussian Markov property this ordered three-block
law is exactly \eqref{chainmeasure}. Equivalently, conditional
expectation of the third block onto the first two is $P_qk$ of
the middle field, leaving
$\int(P_qf)(b)g(b)(P_qk)(b)\,\mu_q(db)$.
Polynomial approximation extends the identity to the stated class,
since h and U are bounded at fixed q.
\end{proof}

This is an application of the established q-Wick rule and
q-Ornstein--Uhlenbeck transitions \cite{bks,eff}. The two transitions
are an auxiliary spectral representation of this contraction;
they do not introduce two independent physical disorder samples.

\newpage
\section{The exact selected two-time function}
Separate the density as
\[
 \rho=f_q(h)+s\lambda(U^2-1)-s\sqrt\lambda\{h,U\},\qquad
 f_q(h)=1-s+sh^2.
\]
The last term is odd under $U\mapsto-U$, whereas $U(u)U(v)$
is even. It therefore vanishes for this observable, but not for
arbitrary mixed observables. Let $\widetilde C_{1,\eps}$ denote
the energy-tilted q-Mehler correlation from \cite{weighted}.
Lemma~\ref{four} gives the exact identity
\begin{align}
 S_\eps(u,v)&:=\sigma(U(u)U(v))
 =\widetilde C_{1,\eps}(u-v)+s\lambda(1+q)B_\eps(u,v),\label{exact}\\
 B_\eps(u,v)&=\int
 e^{i[u(a-b)+v(b-c)]/\sqrt\eps}\,d\Pi_q(a,b,c).\label{B}
\end{align}
The trace correlation is $C_{1,\eps}(u-v)$.
In the first term of \eqref{exact}, cyclicity and $[h,f_q(h)]=0$
combine the exterior phases; this reduction would be invalid
for the full density.

The extra kernel is positive in its two time indices. Indeed
\[
 A_u(b)=e^{-iub/\sqrt\eps}P_q(e^{iu\,\cdot/\sqrt\eps})(b),
 \qquad
 B_\eps(u,v)=\int A_u(b)\overline{A_v(b)}\,\mu_q(db).
\]
In particular $|A_u|\le1$ and $|B_\eps|\le1$. Spectral sign reversal
and chain reversal imply reality and symmetry in u,v.
The semigroup law gives two useful exact reductions:
\begin{equation}\label{slices}
 B_\eps(u,0)=C_{1,\eps}(u),\qquad
 B_\eps(u,u)=C_{2,\eps}(u).
\end{equation}
Consequently
\begin{equation}\label{variance}
 S_\eps(u,u)=1+s\lambda(1+q)C_{2,\eps}(u).
\end{equation}
At fixed $0<q<1$, $C_{2,\eps}(u)\to0$ as $|u|\to\infty$,
by the absolutely continuous gap measure and the
Riemann--Lebesgue lemma. The selected variance therefore
approaches one from its initial value $1+s\lambda(1+q)$.
This fixed-q statement is separate from the compact-time
small-$\eps$ expansion.

The one-point function vanishes. Only the anticommutator term
could contribute, and its two-energy integral contains
$(a+b)e^{iu(b-a)/\sqrt\eps}K_q(a,b)\mu_q(da)\mu_q(db)$.
Exchange of a,b removes its sine part, and simultaneous sign
reversal removes its cosine part. Thus \eqref{variance} is a
variance change, not the square of a nonzero mean.

\newpage
\section{A uniform three-energy estimate}
Use angular energies $e_\theta=2\cos\theta/\sqrt{1-e^{-\eps}}$.
For the middle angle and the two successive gaps write
\[
 \theta_1=\pi/2+\sqrt\eps x/2,\quad
 \theta_0=\theta_1+\eps z_1,\quad
 \theta_2=\theta_1-\eps z_2.
\]
Let $m_\eps(x)$ be the middle marginal density and
$p_\eps(y,z)=p_{1,\eps}(y,z)$ the two-energy density in
\cite{weighted}. Changing variables in \eqref{chainmeasure} gives
\begin{equation}\label{density3}
 j_\eps(x,z_1,z_2)=
 \frac{p_\eps(x+\sqrt\eps z_1,z_1)
       p_\eps(x-\sqrt\eps z_2,z_2)}{m_\eps(x)}
\end{equation}
on the angular domain, extended by zero elsewhere. The Jacobians
in this formula are one for each pair relative to $(x,z_i)$.

Let $\varphi$ be the standard Gaussian density and
$f(z)=2z/\sinh(\pi z)$ with its removable value at zero.
The exact modular marginal is
\[
 m_\eps(x)=\varphi(x)e^{\eps/8}\cos(\sqrt\eps x/2)
 F_\eps(\pi+\sqrt\eps x)Q_\eps^{\rm im}(\pi+\sqrt\eps x).
\]
Here $F_\eps(t)=1-e^{-2\pi\min(t,2\pi-t)/\eps}$ and the image
quotient obeys $|Q^{\rm im}-1|\le4e^{-2\pi^2/\eps}$.
In particular $m_\eps\le C\varphi$.

\begin{lemma}\label{L1}
For every fixed polynomial weight w in $|x|,|z_1|,|z_2|$,
\[
 \int w\,\left|j_\eps-\varphi(x)f(z_1)f(z_2)
       [1-\sqrt\eps x(z_1-z_2)]\right|=O_w(\eps).
\]
\end{lemma}
\begin{proof}
Set $L=\log(1/\eps)$ and restrict first to
$|x|\le8\sqrt L$, $|z_i|\le32L/\pi$. This rectangle lies inside
the angular domain for small eps. The pair expansion of
\cite{weighted} and the marginal above give a relative
$1+O(\eps P)$ factor, with P a fixed polynomial; the Gaussian
ratio is exactly
\[
 \frac{\varphi(x+\sqrt\eps z_1)\varphi(x-\sqrt\eps z_2)}
      {\varphi(x)^2}
 =e^{-\sqrt\eps x(z_1-z_2)-\eps(z_1^2+z_2^2)/2}.
\]
Its remainder after the displayed linear term is bounded by
$C\eps P$ on the rectangle. Integration has no logarithmic loss.
Outside, both pair marginals have the global bound
$p_\eps\le C\varphi f$ from \cite{weighted}; each gap marginal
thus has exponential tails and every fixed moment uniformly.
The middle marginal has Gaussian tails. The union tail probability
is $O(\eps^{32})$ times powers of L. Holder gives uniform joint
polynomial moments from these marginals; Cauchy--Schwarz then
bounds weighted joint tails by $O(\eps^{16})$ times powers of L.
The approximating density has the same harmless tails.
This proves the lemma without dividing by an uncontrolled
small edge marginal.
\end{proof}

\newpage
\section{Time expansion and the surviving neutral memory}
For each adjacent pair let
$\omega_1=(e_{\theta_0}-e_{\theta_1})/\sqrt\eps$ and
$\omega_2=(e_{\theta_1}-e_{\theta_2})/\sqrt\eps$.
The exact trigonometric expressions, with
$y_1=x+\sqrt\eps z_1$, $y_2=x-\sqrt\eps z_2$, are
\[
 \omega_i=
 \frac{-4\cos(\sqrt\eps y_i/2)\sin(\eps z_i/2)}
      {\sqrt{\eps(1-e^{-\eps})}}.
\]
On the entire domain, $|\omega_i|\le C|z_i|$ and
$|\omega_i+2z_i|\le C\eps P(|x|,|z_i|)$.
These follow from real sine/cosine remainders, allowing polynomial
growth. Thus the phase and any fixed number of its time derivatives
may be replaced by their $(-2z_1,-2z_2)$ counterparts at
$O_{S,k}(\eps)$ cost in Lemma~\ref{L1}. The apparent
$\sqrt\eps$ term integrates to zero because it is odd in x.
Using $\int f(z)e^{-2iuz}dz=g(u)$ gives
\begin{equation}\label{Blimit}
 B_\eps(u,v)=g(u)g(v)+O_{C^k([-S,S]^2)}(\eps).
\end{equation}

The previous weighted theorem, now used at d=1, gives
$C_{1,\eps}=g-\eps R/2+O_{C^k}(\eps^2)$ and
$\widetilde C_{1,\eps}=g+\eps R/2+O_{C^k}(\eps^2)$.
Since $s\lambda(1+q)=\eps/2+O(\eps^2)$,
\eqref{exact} and \eqref{Blimit} prove Theorem~\ref{main}.
The new response relative to the trace is
\[
 \frac{S_\eps(u,v)-C_{1,\eps}(u-v)}{\eps}
   =R(u-v)+\tfrac12g(u)g(v)+O_{C^k}(\eps).
\]
The leading neutral addition is a positive rank-one covariance
kernel. That describes its time dependence; it is not a new
particle or a spatial mode identification.

For the equal mixture $\sigma_{\rm mix}=(\tr+\sigma)/2$,
\begin{equation}\label{mix}
 \sigma_{\rm mix}(U(u)U(v))
 =g(u-v)+\tfrac{\eps}{4}g(u)g(v)+O_{C^k}(\eps^2).
\end{equation}
The curvature cancellation of \cite{weighted} therefore does
not restore all two-point functions to their leading form.
More generally, a constant convex mixture with selected weight a
has curvature correction $\eps(a-\tfrac12)R_2$, requiring
$a=1/2$ to cancel it for all times; \eqref{mix} then remains.

Normalizing by the initial selected variance does not remove
this memory either:
\[
 \frac{S_\eps(u,u)}{S_\eps(0,0)}
 =1+\tfrac{\eps}{2}[\sech^4u-1]+O(\eps^2).
\]
Thus the effect is not solely a choice of initial amplitude.

\newpage
\section{An observable that sees the odd mixed source}
The even two-point function alone does not exhaust the source.
Consider the Hermitian energy--slope observable
\[
 A(u)=\{h,U(u)\}.
\]
The microscopic counterpart is the anticommutator of the
normalized Hamiltonian and actual exact-pinched slope; the
fixed-degree preparation interface of \cite{actual} applies.
Its trace and energy-only-tilted expectations vanish by U parity,
but its selected mean does not:
\begin{align}
 \sigma(A(u))&=-s\sqrt\lambda\,H_\eps(u),\label{odd}\\
 H_\eps(u)&=\tr(A(0)A(u))
 =\int(a+b)^2e^{iu(b-a)/\sqrt\eps}
               K_q(a,b)\mu_q(da)\mu_q(db).\nonumber
\end{align}
This follows by keeping $-s\sqrt\lambda\{h,U\}$ in
\eqref{density}; the other terms are odd in U after insertion.
In particular $H_\eps(0)=2(1+q)$.

In the center/gap coordinates of \cite{weighted}, the weight is
exactly
\[
 (e_\theta+e_\phi)^2=
 \frac{16\sin^2(\sqrt\eps y/2)\cos^2(\eps z/2)}
      {1-e^{-\eps}}
 =4y^2+O(\eps P(|y|,|z|)).
\]
It is globally bounded by $Cy^2$. The weighted density remainder
and phase estimate therefore give, with fixed compact-time
derivatives,
\[
 H_\eps(u)=4g(u)+O_{C^k}(\eps),\qquad
 \boxed{\ \sigma(A(u))=-2\sqrt\eps\,\sech^2u
                         +O_{C^k}(\eps^{3/2})\ }.
\]
This is the odd part of the same preparation memory. It is absent
from every energy-only density, not merely from a specific
temperature choice. For the equal mixture its coefficient is
halved and remains nonzero.

An additional independent selected-transfer check verifies 90
polynomial energy--slope identities; 20 physical time points
check the formula, with spectral refinement difference below
$4.45\,10^{-15}$. The truncation argument on the following page
also applies to H: replace the tilted-pair factor F by
$16/(1-q)$ and multiply by $s\sqrt\lambda$ for the mean.
All saved science is kept separate from the analytic remainder
argument.

\newpage
\section{Finite checks and numerical error separation}
The new exact checker verifies \eqref{fouridentity} for every
three-block monomial of total degree at most eight, at
$q=1/4,2/3$ (330 identities). A separate selected open-chord
transfer evaluates $M U\text{-word}\,M^\dagger$ for every
three-block degree at most six, at two rational $(\gamma,\lambda)$
points (168 identities). These are formal algebra checks; the
rational parameters need not satisfy $\gamma=e^{-\lambda}$.

For physical q=0.25,0.64, seven time pairs compare a q-Hermite
Jacobi spectral implementation with direct angular Pochhammer
products (160/224 nodes, 128/192 product terms).
The maximum difference is $3.45\,10^{-14}$; quadrature refinement
is below $3.05\,10^{-14}$. Four small deformations
$\eps=0.16,0.08,0.04,0.02$ use 28 time pairs and Jacobi
dimensions $L=\lceil12/\eps\rceil+64$ and $L+64$.
The largest refinement difference is $1.27\,10^{-15}$;
the maximum phase-sum condition is 1.118.
Exact slices \eqref{slices}--\eqref{variance} are also checked.

An independent saved-output bound calculation controls truncation
in exact arithmetic, rather than treating refinement as a proof.
In the L-node Gauss q-Hermite rule, take
\[
 m=\lfloor(2L-1)/3\rfloor,\quad
 z=\frac{2\max(|u|,|v|,|u-v|)}{\sqrt{\eps(1-q)}},\quad
 \eta=\frac{z^{m+1}}{(m+1)!}.
\]
Real exponential Taylor polynomials have uniform error at most
eta. The q-Mehler map contracts $L^2(\mu_q)$; its finite matrix
is orthogonally conjugate to $\operatorname{diag}(q^n)$ and also
contracts the Gauss-weighted $L^2$ norm. For degree at most m
the two maps agree at Gauss nodes. The product of three
transformed polynomials has degree at most $3m\le2L-1$,
so its quadrature is exact.

Replacing the two exterior exponentials costs at most
$2\eta+\eta^2$ in either representation; replacing the middle
one costs at most $\eta(1+\eta)^2$. Hence
\[
 |B-B_L|\le6\eta+6\eta^2+2\eta^3.
\]
For the tilted pair, $m+2<L$, $2m+2\le2L-1$ and
$F=1-\sqrt q+4\sqrt q/(1-q)$ give
$|\widetilde C-\widetilde C_L|\le F(4\eta+2\eta^2)$.
Combining these bounds with \eqref{exact}, the largest refined
selected truncation bound is $1.44\,10^{-35}$.
Roundoff and angular quadrature discretization are not included
in that analytic bound. Product tails are recorded separately.

The largest sampled selected second remainder increases from
0.454 to 0.494 as eps decreases; this finite coefficient is
consistent with, but does not prove, the error order.
The uniform statement follows from the density proof above.

\newpage
\section{What this adds to the bilocal program}
The actual neutral slope is the first coefficient of the
neutral memory series entering the reconstructed bilocal.
The exact identity \eqref{exact} supplies its previously missing
two-time preparation response in the same microscopic state
used for the curvature. It also produces a direct obstruction
to selecting a universal leading state by the earlier equal
mixture alone.

The next task is to retain the higher neutral residuals $N_j$
and their common-disorder charged cross terms inside the
complete bilocal filter \cite{bilocal}. The odd part
$-\sqrt\lambda\{h,U\}$ is now detected by \eqref{odd}, and must
be retained when propagating the full response.
No conclusion about all bilocal words, spatial generators,
particle masses, finite-G scattering or arbitrary selected
interlacings follows from this single neutral coefficient.

All successful previous scientific outputs are reused read-only.
The companion evidence file records source and output hashes,
dependency versions, both successful LaTeX passes and the
eight-page visual check. This stage is local; no new Git or
external issue reply is performed.

\begin{thebibliography}{9}\small
\bibitem{bks} M. Bo\.zejko, B. K\"ummerer and R. Speicher,
\emph{q-Gaussian processes: non-commutative and classical aspects},
\href{https://arxiv.org/html/funct-an/9604010}{arXiv:funct-an/9604010},
Theorems 2.11, 2.13, 4.4, 4.6 and Corollary 4.5:
conditional expectations and the classical q-Mehler Markov law.
\bibitem{eff} E. G. Effros and M. Popa,
\emph{Feynman diagrams and Wick products associated with q-Fock space},
\href{https://pmc.ncbi.nlm.nih.gov/articles/PMC166362/}{PNAS 100 (2003)},
Corollary 2.1 and Theorem 3.2. The q-Wick machinery is prior work.
\bibitem{selected} A. Day and Chengheng, \emph{Single-selected
covariant density and curvature dynamics}, 8 November 2026;
\texttt{manuscript/susskind-slater-selected-state/main.tex}.
\bibitem{actual} A. Day and Chengheng, \emph{Exact energy
projection in a single Slater excitation}, 8 November 2026;
\texttt{manuscript/susskind-slater-orbit-density/main.tex};
with the preceding neutral-slope normalization.
\bibitem{weighted} A. Day and Chengheng, \emph{A weighted
q-Mehler limit for selected preparations}, 8 November 2026;
\texttt{manuscript/susskind-weighted-mehler-limit/main.tex}.
Its fixed-integer-d theorem is used at d=1 and d=2.
\bibitem{bilocal} A. Day and Chengheng, \emph{Fixed-charge
bilocals and selected preparations}, 8 November 2026;
\texttt{manuscript/susskind-sector-bilocal/main.tex}.
\end{thebibliography}
\end{document}
